\section{Conclusion and future work}
\label{sec:conclusion}

On this setup, from C16 to C32 at the default budget, added delay went into per-step
time while steps filled toward the budget's cap; with the smallest budget it went into
the queue. The budget decides how prompt work is cut into steps and so trades median
against tail latency, and concurrency decides how many steps carry prompt work. A
step's cost follows its tokens and their context closely enough that a four-term model
fit on one-request data predicts the C8 tail, but that result is post hoc, and the
pre-registered test of the token-count relation failed at C8. Four directions follow:
\begin{itemize}
  \item a pre-registered test of the step-cost model in the style of Vidur
    \citep{agrawal2024vidur}: predict ITL p50 and p99 at C32 and other budgets before
    measuring them, and identify the decode term with multi-request data;
  \item open-loop load with varied lengths, to see whether the per-step account and
    packing hold when queues can grow and requests do not arrive in groups;
  \item larger models and multi-GPU serving, where communication adds a third
    component to step cost;
  \item measurement-driven backend selection, which we deferred because single-server
    data cannot evaluate a selection policy without invented rewards.
\end{itemize}
% src: experiments/roadmap.md (Milestone 4 deferral); report next-step sections
